% !TeX root = ../main.tex
\exercise{۲.۶}
\subsection*{صورت مسئله و مدل مورد استفاده}
رابطه \exnum{۲.۲۷} کتاب را برای ولتاژ شروع اشباع، \(V_{\mathrm{dsat}}\)، اثبات کنید
\cite[تمرین \exnum{۲.۶}، صفحه چاپی ۹۶]{weste}.
هدف، اثبات رابطه
\(V_{\mathrm{dsat}}=V_{GT}V_c/(V_{GT}+V_c)\)
در مدل \lr{Velocity Saturation} کتاب است.

در این مدل، قدرمطلق سرعت رانش حامل‌ها \(v\) برحسب قدرمطلق میدان جانبی \(E\) در همان نقطه از کانال به‌صورت قطعه‌ای بیان می‌شود
\cite[روابط \exnum{۲.۲۴} و \exnum{۲.۲۵}، صفحه چاپی ۷۶]{weste}\cite[اسلاید ۹]{nonideal}:
\begin{equation}
 v(E)=
 \begin{cases}
 \dfrac{\mu_{\mathrm{eff}}E}{1+E/E_c},&0\le E<E_c,\\[4pt]
 v_{\mathrm{sat}},&E\ge E_c,
 \end{cases}
 \qquad E_c=\frac{2v_{\mathrm{sat}}}{\mu_{\mathrm{eff}}}.
 \label{eq:velocity-piecewise}
\end{equation}
در اینجا \(\mu_{\mathrm{eff}}\) تحرک مؤثر، \(v_{\mathrm{sat}}\) سرعت اشباع و \(E_c\) میدان بحرانی است.
رابطه \(E_c=2v_{\mathrm{sat}}/\mu_{\mathrm{eff}}\) از پیوستگی سرعت در \(E=E_c\) حاصل می‌شود.
مقدار \(E\) در امتداد کانال تغییر می‌کند. نسبت \(V_{DS}/L\) میدان متوسط را می‌دهد و جایگزین مقدار میدان در هر نقطه نیست.
شاخه نخست رابطه \eqref{eq:velocity-piecewise} نیز فقط برای میدان‌های کمتر از \(E_c\) به کار می‌رود.

ولتاژ اضافه گیت \(V_{GT}=V_{GS}-V_t>0\) و ولتاژ بحرانی \(V_c\) را به‌صورت زیر تعریف می‌کنیم:
\begin{equation}
 V_c=E_cL=\frac{2v_{\mathrm{sat}}L}{\mu_{\mathrm{eff}}}>0,
 \qquad v_{\mathrm{sat}}=\frac{\mu_{\mathrm{eff}}V_c}{2L}.
 \label{eq:velocity-critical-voltage}
\end{equation}
با استفاده از این مدل سرعت، رابطه جریان کتاب دو شاخه دارد
\cite[رابطه \exnum{۲.۲۶}، صفحه چاپی ۷۷]{weste}:
\begin{align}
 I_{\mathrm{lin}}&=\frac{\mu_{\mathrm{eff}}C_{\mathrm{ox}}W}{L}
 \frac{(V_{GT}-V_{DS}/2)V_{DS}}{1+V_{DS}/V_c},
 &&0\le V_{DS}<V_{\mathrm{dsat}},
 \label{eq:velocity-linear}\\
 I_{\mathrm{sat}}&=C_{\mathrm{ox}}W(V_{GT}-V_{\mathrm{dsat}})v_{\mathrm{sat}},
 &&V_{DS}\ge V_{\mathrm{dsat}}.
 \label{eq:velocity-saturation}
\end{align}
کمیت \(q'_{\mathrm{inv,d}}=C_{\mathrm{ox}}W(V_{GT}-V_{\mathrm{dsat}})\)، قدرمطلق بار وارونگی بر واحد طول در انتهای درین است.
واحد آن \(\unit{\coulomb\per\metre}\) است؛ ضرب آن در سرعت \(v_{\mathrm{sat}}\) جریان را با واحد \(\unit{\ampere}\) می‌دهد.
بنابراین این کمیت، بار کل کانال نیست.

\subsection*{پیوستگی جریان و استخراج ولتاژ اشباع}
برای اختصار، \(A=V_{GT}\) و \(s=V_{\mathrm{dsat}}\) می‌گذاریم.
پیوستگی دو شاخه جریان در \(V_{DS}=s\)، با جایگذاری رابطه \eqref{eq:velocity-critical-voltage}، می‌دهد
\begin{equation}
 \frac{\mu_{\mathrm{eff}}C_{\mathrm{ox}}W}{L}
 \frac{(A-s/2)s}{1+s/V_c}
 =C_{\mathrm{ox}}W(A-s)\frac{\mu_{\mathrm{eff}}V_c}{2L}.
 \label{eq:velocity-continuity}
\end{equation}
با حذف عوامل مشترک مثبت و ضرب در \(2(1+s/V_c)\)، رابطه ساده می‌شود:
\begin{align}
 2\left(A-\frac{s}{2}\right)s&=(A-s)(V_c+s),\notag\\
 2As-s^2&=AV_c+As-V_cs-s^2,\notag\\
 (A+V_c)s&=AV_c.
 \label{eq:velocity-algebra}
\end{align}
جمله‌های درجه‌دو از دو طرف حذف شده‌اند؛ چون \(A+V_c>0\)، نتیجه یکتا است:
\begin{equation}
 \boxed{V_{\mathrm{dsat}}=\frac{V_{GT}V_c}{V_{GT}+V_c}
 =V_{GT}\parallel V_c.}
 \label{eq:velocity-result}
\end{equation}
نماد \(a\parallel b=ab/(a+b)\) در اینجا فقط تشابه جبری با فرمول مقاومت‌های موازی را نشان می‌دهد.
برای پارامترهای مثبت، نتیجه \eqref{eq:velocity-result} شرط
\(0<V_{\mathrm{dsat}}<\min(V_{GT},V_c)\) را برآورده می‌کند.

\subsection*{کنترل رفتار حدی و دامنه اعتبار}
با جایگذاری رابطه \eqref{eq:velocity-result} در جریان اشباع داریم
\cite[رابطه \exnum{۲.۲۸}]{weste}
\begin{equation}
 I_{\mathrm{sat}}=C_{\mathrm{ox}}Wv_{\mathrm{sat}}
 \frac{V_{GT}^2}{V_{GT}+V_c}.
 \label{eq:velocity-final-current}
\end{equation}
در حد \(V_c\gg V_{GT}\)، ولتاژ شروع اشباع به \(V_{GT}\) نزدیک می‌شود.
با توجه به رابطه \eqref{eq:velocity-critical-voltage}، جریان \eqref{eq:velocity-final-current} در این حد به
\(I_{\mathrm{sat}}\approx\mu_{\mathrm{eff}}C_{\mathrm{ox}}(W/L)V_{GT}^2/2\)
تبدیل می‌شود و مدل مربعی کانال بلند بازیابی می‌شود.

در حد \(V_{GT}\gg V_c\)، داریم
\(V_{\mathrm{dsat}}\approx V_c\) و
\(I_{\mathrm{sat}}\approx C_{\mathrm{ox}}Wv_{\mathrm{sat}}V_{GT}\).
بنابراین، با ثابت گرفتن پارامترهای مدل، جریان در این حد تقریباً با \(V_{GT}\) متناسب است.
این نتیجه درباره حد مدل است؛ استقلال کامل ولتاژ اشباع از گیت یا اعتبار یکسان آن برای همه فناوری‌ها را ثابت نمی‌کند.
خود \(\mu_{\mathrm{eff}}\) می‌تواند با ولتاژ گیت تغییر کند و در نتیجه \(V_c\) نیز لزوماً ثابت نیست
\cite[صفحه چاپی ۷۷]{weste}.
